\exnumber{02}
\extitle{UI IV - Ex 3 - Design screens}
\exfiles{ex02/Figma file, named Name\_FirstName\_DeliverableName\_Date\_VersionNumber}
\exauthorize{Those named in the exercise}
\makeheaderfiles

\textbf{Instructions:}

Based on the wireframes and your design system library, take your mid-fidelity prototype to a high-fidelity level!

Select 5 screens that will show the user flow of a customer who:

\begin{itemize}
\item arrives on the homepage
\item browses the content
\item checks one product page
\item adds the product to the cart
\item checks out
\end{itemize}

Use your components, call to actions and visuals to complete the wireframes.

\textbf{Deliverables:}

\begin{itemize}
\item 5 high-fidelity screens from Figma
\item Upload the Figma file on Github, labeled as follows: "Name\_FirstName\_DeliverableName\_Date\_VersionNumber". For example, it can be labeled as "Doe\_John\_MidFiPrototype\_05242024\_V1".
\item Add a title within the file.
\item Pay attention to the global aspect of the file. It must be clear, simple and easy to read. You can get inspiration on canva.com but don't overload your design with too much details!
\end{itemize}

\textbf{Resources:}

\begin{itemize}
\item Unsplash \url{https://unsplash.com/} for visuals
\item From low-fidelity to high-fidelity \url{https://mockitt.wondershare.com/prototyping/low-fidelity-vs-high-fidelity.html}
\end{itemize}

